%% ═══════════════════════════════════════════════════════════════════════════════
\section{Transfer Mode: Verifying Robustness Across Networks}
\label{sec:transfer}
%% ═══════════════════════════════════════════════════════════════════════════════

In this section, we present our main contribution: the transfer-mode formulation. We begin by motivating the problem, then describe our three-network MIP encoding, state the transfer MIP, and prove its correctness.

\subsection{Motivation}

Training a more accurate network $N_2$ from $N_1$ (e.g., via distillation or fine-tuning) often increases prediction confidence on correctly classified inputs. A natural question is: \emph{by how much can $N_2$'s confidence exceed $N_1$'s in the worst case, in the regime where $N_1$ is already certifiably robust but $N_2$ can still be fooled?} The answer to this question determines whether $N_1$'s robustness guarantee transfers to $N_2$.

\subsection{Three-Network Encoding}

To answer this question, the transfer mode encodes \textbf{three} forward passes simultaneously in a single MIP: $N_1(x)$, $N_2(x)$, and $N_2(\xp)$.

\begin{center}
\begin{tikzpicture}[
  box/.style={draw,rounded corners,minimum width=2.4cm,minimum height=0.9cm,
               text centered,fill=blue!10},
  arr/.style={-Stealth,thick}]
  \node[box] (n1c) {$N_1(x)$};
  \node[box,right=2cm of n1c] (n2c) {$N_2(x)$};
  \node[box,below=1.2cm of n2c] (n2p) {$N_2(\xp)$};
  \draw[arr] (n1c) -- node[above,font=\small]{same $x$} (n2c);
  \draw[arr,dashed] (n2c) -- node[right,font=\small]{$\xp\in\mI_\ep(x)$} (n2p);
\end{tikzpicture}
\end{center}

Each forward pass uses a unique variable-name prefix to distinguish the three copies. The clean input $x$ is shared between $N_1(x)$ and $N_2(x)$, and the perturbed input $\xp$ is linked to $x$ via the perturbation constraint. Note that $N_1(\xp)$ is \emph{not} encoded: the transfer MIP only requires $N_1$'s confidence on the clean input to establish the robustness baseline, and $N_2$'s computation on both the clean and perturbed inputs to verify whether the attack succeeds. This saves all of the ReLU binary variables that would be needed for a fourth network pass, significantly reducing the problem size.

\subsection{Transfer MIP Formulation}

Let $c_1 = \conf{N_1}{x}{c}$, $c_2 = \conf{N_2}{x}{c}$, and $c_2' = \conf{N_2}{\xp}{c}$.
The transfer MIP is:

\begin{equation}
  \label{eq:transfer}
  \dD \;=\; \max_{x,\,\xp,\,e}\; (c_2 - c_1)
  \quad\text{s.t.}\quad
  \begin{cases}
    \xp \in \mI_\ep(x) & \text{(perturbation constraint)} \\
    c_1 \;\ge\; \dO + \epsilon_0 & \text{(T1: $N_1$ is certifiably robust)} \\
    c_2 - c_1 \;\ge\; 0 & \text{(T2: $N_2$ at least as confident as $N_1$)} \\
    \text{attack constraint on } N_2(\xp) & \text{(T3: $N_2$ can be fooled)}
  \end{cases}
\end{equation}

where $\epsilon_0 = 10^{-3}$ is a strict-inequality guard. In the formal notation of Section~\ref{sec:notation:formal}, this is exactly the transfer MIP of Definition~\ref{def:ddiff}, with $\dD = \ddiff$.

The constraint T1 restricts the search to inputs where $N_1$ is certifiably robust (its confidence exceeds the standard-mode bound $\dO$). Constraint T2 focuses on the regime where $N_2$ is more confident than $N_1$. Constraint T3 requires that a perturbation exists that fools $N_2$, making the transfer problem non-trivial. The encoding of T3 depends on the chosen attack mode (Section~\ref{sec:modes}).

\subsection{Correctness of the Transfer Guarantee}

We now prove that the transfer MIP provides a sound robustness certificate. The key insight is that any input violating the transfer guarantee would be a feasible solution to the MIP with objective value exceeding $\dD$, contradicting its optimality.

\begin{theorem}[Transfer Robustness]
  \label{thm:transfer}
  Let $\dO$ be the standard-mode value for $N_1$ on perturbation $(\ep, c \to t)$. Let $\dD$ be the optimal value of the transfer MIP~\eqref{eq:transfer}. Then for any input $x$:
  \[
    \conf{N_1}{x}{c} \;\ge\; \dO
    \;\;\text{and}\;\;
    \conf{N_2}{x}{c} - \conf{N_1}{x}{c} \;>\; \dD
    \;\;\implies\;\; x \text{ is } (\ep,c\to t)\text{-robust for } N_2.
  \]
\end{theorem}

\begin{proof}
  Suppose for contradiction that $x$ is \emph{not} $(\ep,c\to t)$-robust for $N_2$. Then there exists $\xp \in \mI_\ep(x)$ such that $N_2(\xp) = t$, meaning the attack constraint T3 holds.

  Since $\conf{N_1}{x}{c} \ge \dO$, constraint T1 is satisfied (with the $\epsilon_0$ guard). Since $\conf{N_2}{x}{c} - \conf{N_1}{x}{c} > \dD \geq 0$, constraint T2 is also satisfied.

  Therefore $(x, \xp)$ is a feasible solution to the transfer MIP~\eqref{eq:transfer} with objective value $c_2 - c_1 > \dD$, contradicting the fact that $\dD$ is the maximum feasible objective value.
\end{proof}

\begin{corollary}[Single-threshold inference]
  \label{cor:single}
  Define $\Delta := \dO + \dD$. If $\conf{N_1}{x}{c} \ge \dO$ and $\conf{N_2}{x}{c} > \Delta$, then $x$ is $(\ep, c \to t)$-robust for $N_2$.
\end{corollary}

This corollary shows that at inference time, the transfer guarantee reduces to checking two confidence thresholds, requiring only two forward passes and no additional optimization.

\begin{remark}
  Checking $\conf{N_2}{x}{c} > \Delta$ alone (without verifying $N_1$'s threshold) is \emph{unsound}: constraint T1 would not be guaranteed, making the transfer MIP potentially infeasible and $\dD$ meaningless. Both conditions must be verified.
\end{remark}

\subsection{Relationship to the Formal Quantities}

The transfer MIP~\eqref{eq:transfer} computes $\dD = \ddiff$, which satisfies $\dN{2} \geq \dN{1} + \ddiff$ (Theorem~\ref{thm:correct} in Section~\ref{sec:formal}). This is a \emph{lower bound}: knowing $\dD$ certifies that $N_2$'s standard-mode value is at least $\dO + \dD$, but $N_2$'s actual standard-mode value may be strictly larger. The precise gap and a counterexample are analyzed in Section~\ref{sec:formal}.
